\documentclass[a4paper,11pt]{ltjsarticle}
\usepackage{amsmath, amssymb}
\usepackage{bm}
\usepackage{geometry}
\geometry{margin=25mm}

\title{チーム編成ログからのキャラクター相性抽出と3人目推薦のモデル化}
\author{}
\date{}

\begin{document}
\maketitle

\section{「相性」を2種類の $N \times N$ 行列（ヒートマップ）に落とす}

3人組のままだと組み合わせ爆発（${}_{80}\mathrm{C}_3 \approx 8\text{万通り}$）で1マスあたりの試合数が疎になるため、まずはペア（2体関係）に分解して集計する。

\begin{itemize}
  \item \textbf{味方シナジー行列（Synergy）}：同じチームにキャラクター $A$ と $B$ が同時に編成された試合の勝率。
  \item \textbf{対面カウンター行列（Matchup）}：自チームに $A$、相手チームに $B$ がいた試合の $A$ 側の勝率。
\end{itemize}

\subsection*{単体勝率からのリフト値（残差）による定義}

単なる生勝率でヒートマップを描くと、「そもそも単体スペックが高くて勝率が高い強キャラ」の行・列が全体的に高くなるだけで、純粋な「組み合わせの相性」が見えなくなる。そこで、単体勝率からの期待値との差分（リフト値・残差）をとって行列要素を定義する。

\begin{equation}
  \mathrm{Synergy}(A, B) = P(\mathrm{Win} \mid A, B \in \mathrm{Team}) - \frac{P(\mathrm{Win} \mid A) + P(\mathrm{Win} \mid B)}{2}
\end{equation}

同様に対面カウンター行列についても、単体勝率の期待値からの差分として定義する。

\begin{equation}
  \mathrm{Matchup}(A, B) = P(\mathrm{Win}_A \mid A \in \mathrm{Team}, B \in \mathrm{Enemy}) - \left( 0.5 + \frac{P(\mathrm{Win} \mid A) - P(\mathrm{Win} \mid B)}{2} \right)
\end{equation}

※ 味方シナジー行列と対面カウンター行列は意味が逆であるため、足し合わせずに別々の行列として保持する。

\section{共起・シナジーネットワーク図（3人組の発見）}

上記の「シナジー値が高いペア」だけをエッジで結んだ無向グラフ（ばねモデルによる自動レイアウト）を描画する。

\begin{itemize}
  \item 相性の良いペア同士が引き合うように配置されるため、結果として「環境で強い3人組（トリオ）」が三角形のクラスター（コミュニティ）として視覚的に浮き上がる，ことを期待。
\end{itemize}

\section{ロジスティック回帰（拡張Bradley-Terryモデル）によるパラメータ推定}

単純集計では「対戦数が少ないマイナーな組み合わせ」の勝率が極端（0\%や100\%）に振れてノイズになる。これを抑えるため、勝敗 $y \in \{0, 1\}$ を以下の特徴量でL2（ユークリッド距離）正則化付きロジスティック回帰にかけ、パラメータを推定する。

\begin{itemize}
  \item 自チームのキャラクター $i$ の存在フラグ（$x_i = +1$）、敵チームのキャラクター $j$ の存在フラグ（$x_j = -1$） \\
  $\rightarrow$ 回帰係数 $\beta_i$ が\textbf{「キャラクター単体の強さ」}を表す。
  \item 自チーム内のペア $(i, j)$ の共起フラグ（$z_{ij} = +1$）、敵チーム内のペア $(k, l)$ の共起フラグ（$z_{kl} = -1$） \\
  $\rightarrow$ 回帰係数 $\gamma_{ij}$ が\textbf{「単体の強さを差し引いた純粋な味方シナジー（相性値）」}を表す。
  \item 自チームの $i$ と敵チームの $j$ の対面フラグ（$w_{ij} = +1$） \\
  $\rightarrow$ 回帰係数 $\delta_{ij}$ が\textbf{「対面カウンターの相性値」}を表す。
\end{itemize}

\section{「2人決まった状態での3人目選び」のモデル化}

組み合わせの対称性・交換可能性を仮定する。
ログに選択順序（ピック順）が記録されていなくても、任意の3人組 $\{A, B, C\}$ の試合ログは、「$A$ と $B$ が先に決まっていて、3人目に $C$ を選んだ状況」のサンプルとしてそのまま扱うことができる。

すでに確定している味方2体 $\{A, B\}$ および相手チーム $\mathrm{Enemy}$ に対し、3人目の候補キャラクター $C$ を加えたときの評価スコアを以下のように定義する。

\begin{equation}
  \mathrm{Score}(C \mid A, B, \mathrm{Enemy}) = \underbrace{\beta_C}_{C\text{単体の強さ}} + \underbrace{\gamma_{AC} + \gamma_{BC}}_{\text{味方2人とのシナジー}} + \underbrace{\sum_{E \in \mathrm{Enemy}} \delta_{CE}}_{\text{敵チームへのカウンター}}
\end{equation}

これにより、味方に $A$ と $B$ が確定し、敵に $\mathrm{Enemy}$ が見えている状況において、どのキャラクター $C$ を選べば勝率が最大になるかをランキング形式で算出できる。

\section{視覚化のアイデア（3人目の推薦をどう見せるか）}

\begin{itemize}
  \item \textbf{条件付き勝率の棒グラフ（またはランキング表）}：\\
  代表的なペア（例：タンク＋ヒーラーなど）をいくつか固定し、「この2人が決まっているとき、3人目に誰を入れると勝率が最も上昇するか（あるいは低下するか）」をソートして並べる。
  \item \textbf{役割（ロール）への集約}：\\
  全キャラクターのままだと組み合わせが細かすぎるため、公式のクラス分類（タンク、アサシン、サポート、スナイパー等）に集約し、「$\{\text{タンク}, \text{サポート}\}$ の2人が決まったとき、3人目はどのロールを選ぶべきか」という $N_{\mathrm{role}}^2 \times N_{\mathrm{role}}$ の小さな表にまとめる。これによりデータ不足（スパース性）が解消され、人間が理解できる戦略（定石）として整理しやすくなる。
\end{itemize}

\end{document}
